\chapter{Mô tả khảo sát và đặt vấn đề}

\section{Bối cảnh}
Việc lựa chọn ngành đại học phù hợp năng lực học tập ở bậc phổ thông là một quyết định có tác động dài hạn đối với học sinh và gia đình. Ở Việt Nam, điểm các môn thi tốt nghiệp THPT (thang mười theo chuẩn Bộ GD\&ĐT) là tín hiệu thống nhất, dễ thu thập và có tính pháp lý cao, trong khi các hệ thống tư vấn thuần học máy thường thiếu giải thích trực quan, còn hệ thống dựa luật cứng nhắc khó bắt được cấu trúc phức tạp của dữ liệu thực.

Dự án \texttt{KBS} (phiên bản 3.0) trong workspace nghiên cứu xây dựng một \textbf{hệ gợi ý lai ghép}: nhánh \textbf{học máy} học quan hệ giữa sáu điểm môn theo từng khối thi và nhãn ngành được gán tự động; nhánh \textbf{hệ cơ sở tri thức (KBS)} mã hóa kinh nghiệm chuyên gia dưới dạng luật có ngưỡng; tầng \textbf{lai ghép (fusion)} và cơ chế \textbf{VETO} điều phối hai nguồn khi chúng mâu thuẫn. Giao diện người dùng triển khai bằng Streamlit (\texttt{app.py}), phù hợp trình diễn và thử nghiệm nhanh.

\section{Vấn đề nghiên cứu}
Các câu hỏi trung tâm của báo cáo gồm:
\begin{enumerate}[leftmargin=*]
  \item Làm thế nào xây dựng tập dữ liệu huấn luyện \textbf{có thể lặp lại} từ file điểm THPT 2024, phản ánh đúng hai khối KHTN và KHXH (sáu môn mỗi khối, không dùng Tin học)?
  \item Cách biểu diễn và thực thi luật KBS trên JSON sao cho \textbf{dễ bảo trì}, có giải thích tiếng Việt và xử lý xung đột khi nhiều luật cùng khớp?
  \item Thuật toán Random Forest với siêu tham số cố định trong \texttt{kbs/config.py} có đạt độ phân tách tốt trên nhãn heuristic hay không, và mức đóng góp từng đặc trưng ra sao?
  \item Công thức lai ghép $0{,}6\,\text{ML} + 0{,}4\,\text{KBS}$ cùng VETO ảnh hưởng thế nào tới \textbf{độ khớp nhãn} so với RF thuần, và ý nghĩa thực tiễn khi nhãn chỉ là xấp xỉ?
\end{enumerate}

\section{Mục tiêu}
\begin{itemize}[leftmargin=*]
  \item \textbf{Mục tiêu chức năng:} Cho vector sáu điểm theo đúng thứ tự đặc trưng của khối, hệ thống xuất bảng xế hạng các ngành trong khối kèm điểm lai, điểm ML, điểm KBS và văn bản giải thích ngắn.
  \item \textbf{Mục tiêu kỹ thuật:} Tách module cấu hình (\texttt{kbs/config.py}), engine luật (\texttt{kbs/knowledge\_rules.py}), fusion (\texttt{kbs/hybrid\_fusion.py}), script dữ liệu và huấn luyện (\texttt{scripts/}), đảm bảo có thể tái tạo pipeline từ dữ liệu gốc tới artifact \texttt{.pkl}.
  \item \textbf{Mục tiêu đánh giá:} Báo cáo định lượng độ chính xác/F1 trên tập kiểm tra và ma trận nhầm lẫn cho từng khối; bổ sung đánh giá hệ lai trên mẫu con có kích thước nêu rõ để tiết kiệm chi phí tính toán khi suy luận từng mẫu phải lặp qua tất cả ngành.
\end{itemize}

\section{Phạm vi và giới hạn}
\textbf{Phạm vi:} Dữ liệu một năm (\texttt{data/diem\_thi\_thpt\_2024.csv}); hai khối; chín ngành tổng hợp trong ánh xạ \texttt{NGANH\_HOC\_MAP} nhưng mỗi khối chỉ xét năm hoặc bốn nhãn tương ứng; ngôn ngữ giao diện và giải thích tiếng Việt.

\textbf{Giới hạn:} Nhãn \texttt{nganh\_hoc} được tạo bởi heuristic trong \texttt{scripts/create\_data.py}, không thay thế khảo sát định hướng nghề nghiệp với chuyên gia tâm lý giáo dục; độ chính xác học máy cao một phần phản ánh tính nhất quán của quy tắc gán nhãn với vector đặc trưng. Luật KBS chưa qua hội đồng giáo viên phản biện. Hệ thống không dự báo điểm chuẩn hay xét tuyển theo tổ hợp cụ thể của từng trường.

\section{Cấu trúc báo cáo}
Chương 2 trình bày cơ sở lý thuyết và công nghệ. Chương 3 mô tả dữ liệu và pipeline. Chương 4--6 trình KBS, huấn luyện AI và tích hợp lai. Chương 7 báo cáo số liệu đánh giá từ mã nguồn. Chương 8--9 mô tả vận hành cập nhật, kết luận. Phụ lục đính kèm đoạn cấu hình và hướng dẫn biên dịch.

\section{Khảo sát tài liệu và mã nguồn liên quan}
Ngoài mã nguồn, nhóm tài liệu \texttt{docs/*\_v3.md} được dùng làm cơ sở mô tả ngữ nghĩa dữ liệu, luật và khung đánh giá bảy bước \cite{docs_hybrid}. README gốc tóm tắt lệnh chạy nhanh và phụ thuộc Python \cite{docs_dataset}. Các tài liệu này được chuyển hóa thành nội dung giải thích trong các chương kỹ thuật, đồng thời bổ sung số liệu đo trực tiếp trên máy chủ mã (xem chương 7).

\section{Tính cấp thiết của đề tài}
Khối lượng thông tin tuyển sinh lớn trong khi thời gian ra quyết định có hạn tạo áp lực nhận thức (\textit{cognitive load}) cho học sinh. Một hệ trợ giúp \textbf{tổng hợp} điểm thành xế hạng ngành kèm giải thích giúp giảm không gian lựa chọn và tăng tính minh bạch so với chỉ dựa cảm tính. Đồ án đóng góp mã nguồn có thể tái lập và tài liệu đi kèm, phù hợp triển khai tại phòng máy trường hoặc máy tính cá nhân.

\section{Đối tượng và phạm vi nghiên cứu}
\textbf{Đối tượng:} repository \texttt{KBS} phiên bản 3.0, gồm pipeline dữ liệu, hai mô hình RF, engine luật JSON, fusion và ứng dụng Streamlit.

\textbf{Phạm vi ngoài đồ án:} không xây dựng cổng đăng nhập người dùng; không tích hợp CSDL tuyển sinh Bộ; không khảo sát định lượng mức độ hài lòng người dùng thật (có thể là hướng mở rộng sau tốt nghiệp).

\section{Phương pháp nghiên cứu}
\begin{enumerate}[leftmargin=*]
  \item \textbf{Phân tích tĩnh và thiết kế:} đọc mã nguồn, sơ đồ hóa luồng dữ liệu và phụ thuộc module.
  \item \textbf{Thực nghiệm tái lập:} chạy \texttt{create\_data}, \texttt{train\_model}, \texttt{evaluate\_model} trên máy trạm, ghi nhận metric và ma trận nhầm lẫn.
  \item \textbf{Đối chiếu lý thuyết:} so khớp thiết kế với tài liệu học thuật về RF, hệ chuyên gia và lai ghép (chương 2).
\end{enumerate}

\section{Đóng góp dự kiến}
\begin{itemize}[leftmargin=*]
  \item Một ứng dụng đầu cuối (\textit{end-to-end}) minh họa lai ghép KBS--ML trên miền giáo dục.
  \item Phân tích trung thực giới hạn nhãn proxy và ảnh hưởng tới độ đo định lượng.
  \item Đề xuất lộ trình API, đa năm dữ liệu, hiệu đính luật bởi chuyên gia và kiểm thử người dùng.
\end{itemize}
